\section{大模型时代的数据规模}
\href{https://epoch.ai/data/ai-models}{数据规模}
pytorch在执行GPU上的运算时，是异步的，如果需要GPU计算的结果，
需要调用@torch.cuda.synchronize()@来强制同步。
如果需要随机数生成的结果可复现，需要设置随机数种子。
\begin{code}
\begin{verbatim}
def set_seed(seed: int = 37) -> None:
    np.random.seed(seed)
    random.seed(seed)

    torch.manual_seed(seed)
    torch.cuda.manual_seed(seed)
    torch.backends.cudnn.deterministic = True
    torch.backends.cudnn.benchmark = False

    os.environ["PYTHONHASHSEED"] = str(seed)
    print(f"Random seed set to: {seed}")
\end{verbatim}
\end{code}

GPU锁频：
\begin{code}
\begin{verbatim}
nvidia-smi --query-gpu=pstate,clocks.mem,clocks.sm,clocks.gr --format=csv
nvidia-smi --query-supported-clocks=gpu_name,mem,gr --format=csv
sudo nvidia-smi -pm 1
nvidia-smi -ac 8001,2100
\end{verbatim}
\end{code}

CPU锁频：
\begin{code}
\begin{verbatim}
sudo apt install cpufrequtils

cpufreq-info
# 设置最大最小频率
sudo cpufreq-set -r -g performance
sudo cpufreq-set -r -d 2.1Ghz
sudo cpufreq-set -r -u 2.1Ghz
\end{verbatim}
\end{code}

测量程序运行时间
\begin{code}
\begin{verbatim}
time.perf_counter() //单位为s

torch.cuda.synchronize()
start=torch.cuda.Event(enable_timing=True)
end=torch.cuda.Event(enable_timing=True)

start.record()
# 运行代码
end.record()
torch.cuda.synchronize()

print(f"{start.elapsed_time(end)}ms")
\end{verbatim}
\end{code}
